Let \(P^{\mathrm F}\in\mathcal M^{\mathrm F}\), and let \(P\in\mathcal M\) be its observed margin.  Definition \(\texttt{def:full-data-model-class}\) supplies, for this particular member, a standard Borel carrier \(\mathcal U\), an \(\mathcal U\)-valued variable \(U\), jointly Borel maps \(g_x\), and the process \(Y(a)\); it also says that Assumptions \(\texttt{ass:consistency}\), \(\texttt{ass:exchangeability}\), and \(\texttt{ass:full-data-response-continuity}\) hold within this carrier.  Fix \(x\in\mathcal X\).  Assumption \(\texttt{ass:stratum-mass}\) gives
\[
 P^{\mathrm F}(X=x)=P(X=x)=p_x\geq p_{\min}>0,
 \tag{1}
\]
so conditioning on \(X=x\) is well defined.  Assumption \(\texttt{ass:consistency}\) gives one \(P^{\mathrm F}\)-probability-one event on which, simultaneously for every \(a\in[0,1]\),
\[
 Y(a)=g_X(a,U),
 \qquad
 Y=g_X(A,U).
 \tag{2}
\]
The simultaneous form in (2) permits evaluation at the random argument \(A\).  Assumption \(\texttt{ass:exchangeability}\) states \(A\perp\!\!\!\perp U\mid X\).  Because \([0,1]\) and the member's \(\mathcal U\) are standard Borel and (1) gives positive stratum mass, the defining conditional-independence factorization is
\[
 P^{\mathrm F}_{(A,U)\mid X=x}
 =P^{\mathrm F}_{A\mid X=x}\otimes P^{\mathrm F}_{U\mid X=x}.
 \tag{3}
\]
Define, within this same member,
\[
 m_x^{\mathrm F}(a)
 =\int_{\mathcal U}g_x(a,u)\,P^{\mathrm F}(\mathrm du\mid X=x)
 =\mathbb E_{P^{\mathrm F}}[g_x(a,U)\mid X=x].
 \tag{4}
\]
All integrands are Borel measurable and bounded by one, by the packaged-map condition in Definition \(\texttt{def:full-data-model-class}\), so the following product integrals are finite.  For any Borel \(B\subseteq[0,1]\), equations (2)--(4) and Assumption \(\texttt{ass:conditional-density-law}\) imply
\[
\begin{aligned}
 \mathbb E_{P^{\mathrm F}}[Y\mathbf 1\{A\in B\}\mid X=x]
 &=\int_{B\times\mathcal U}g_x(a,u)\,
   \{P^{\mathrm F}_{A\mid X=x}\otimes P^{\mathrm F}_{U\mid X=x}\}
   (\mathrm da,\mathrm du)\\
 &=\int_B\left\{\int_{\mathcal U}g_x(a,u)\,
   P^{\mathrm F}(\mathrm du\mid X=x)\right\}
   P^{\mathrm F}(\mathrm da\mid X=x)\\
 &=\int_Bm_x^{\mathrm F}(a)\pi_x(a)\,\mathrm da.
\end{aligned}
\tag{5}
\]
Thus the asserted conditional integral identity is fiberwise: it uses only the carrier packaged by the chosen law.

We next derive, rather than assume, equality with the declared continuous observed-regression version.  Equations (2) and (4), the triangle inequality, and Assumption \(\texttt{ass:full-data-response-continuity}\) yield, for every \(s\in[0,1]\),
\[
\begin{aligned}
 |m_x^{\mathrm F}(t)-m_x^{\mathrm F}(s)|
 &\leq\mathbb E_{P^{\mathrm F}}
   [|g_x(t,U)-g_x(s,U)|\mid X=x]\\
 &=\mathbb E_{P^{\mathrm F}}
   [|Y(t)-Y(s)|\mid X=x]
 \longrightarrow0
\end{aligned}
\tag{6}
\]
as \(t\to s\).  Hence \(m_x^{\mathrm F}\) is continuous.  Assumption \(\texttt{ass:holder-regression}\) says that \(\mu_x^P\) is continuous and is an almost-everywhere version of \(\mathbb E_P[Y\mid A=a,X=x]\).  Assumption \(\texttt{ass:conditional-density-law}\) therefore gives
\[
 \mathbb E_P[Y\mathbf 1\{A\in B\}\mid X=x]
 =\int_B\mu_x^P(a)\pi_x(a)\,\mathrm da.
 \tag{7}
\]
The observed-margin identity and (5)--(7) imply
\[
 \int_B\{m_x^{\mathrm F}(a)-\mu_x^P(a)\}\pi_x(a)\,\mathrm da=0
 \quad\text{for every Borel }B\subseteq[0,1].
 \tag{8}
\]
Put \(f=m_x^{\mathrm F}-\mu_x^P\).  For every integer \(j\geq1\), apply (8) with \(B=\{f\geq1/j\}\) to obtain
\[
 0=\int_{\{f\geq1/j\}}f(a)\pi_x(a)\,\mathrm da
 \geq \frac1j\int_{\{f\geq1/j\}}\pi_x(a)\,\mathrm da.
 \tag{9}
\]
Applying (8) to \(\{f\leq-1/j\}\) gives the corresponding negative-side conclusion.  Thus \(f=0\) for \(\pi_x(a)\,\mathrm da\)-almost every \(a\).  Assumption \(\texttt{ass:polynomial-thinning}\), with \(c_->0\) and \(\kappa\geq0\), gives
\[
 \pi_x(a)\geq c_-a^\kappa>0
 \quad\text{for Lebesgue-almost every }a\in(0,1].
 \tag{10}
\]
Consequently \(f=0\) Lebesgue-almost everywhere on \((0,1]\).  If \(f(a_0)\neq0\) at any \(a_0\in[0,1]\), continuity of both terms would produce a relative neighborhood of \(a_0\) on which \(|f|\) is bounded away from zero; its intersection with \((0,1]\) has positive Lebesgue measure, contradicting the preceding equality.  Therefore
\[
 m_x^{\mathrm F}(a)=\mu_x^P(a)
 \quad\text{for every }a\in[0,1].
 \tag{11}
\]
In particular, equality at the boundary \(a=0\) follows from continuity and is not imposed by choosing a conditional-expectation version.

Fix \(\delta\in[0,\bar\delta]\).  Definition \(\texttt{def:clamp-policy}\) and the simultaneous identities (2) give on the same probability-one event
\[
\begin{aligned}
 Y\{d_\delta(A)\}
 &=g_X\{\max(A,\delta),U\}\\
 &=g_X(A,U)\mathbf 1\{A>\delta\}
   +g_X(\delta,U)\mathbf 1\{A\leq\delta\}\\
 &=Y\mathbf 1\{A>\delta\}
   +Y(\delta)\mathbf 1\{A\leq\delta\}.
\end{aligned}
\tag{12}
\]
This includes \(A=\delta\), because (2) then gives \(Y=Y(\delta)\).  By (3), (4), (11), and Proposition \(\texttt{prop:pushforward-setup}\),
\[
\begin{aligned}
 \mathbb E_{P^{\mathrm F}}[Y(\delta)\mathbf 1\{A\leq\delta\}\mid X=x]
 &=m_x^{\mathrm F}(\delta)P^{\mathrm F}(A\leq\delta\mid X=x)\\
 &=\mu_x^P(\delta)q_x(\delta).
\end{aligned}
\tag{13}
\]
Taking expectations in (12), using (13), using the observed-margin identity for the retained natural-course term, and summing over the finite set \(\mathcal X\), we obtain
\[
\begin{aligned}
 \psi_\delta(P^{\mathrm F})
 &=\mathbb E_{P^{\mathrm F}}[Y\{d_\delta(A)\}]\\
 &=\mathbb E_P[Y\mathbf 1\{A>\delta\}]
   +\sum_{x\in\mathcal X}p_xq_x(\delta)\mu_x^P(\delta)\\
 &=\theta_\delta(P).
\end{aligned}
\tag{14}
\]
The first and last equalities are the definitions of \(\psi_\delta\) and \(\theta_\delta\).  Proposition \(\texttt{prop:policy-support}\) supplies the remaining support assertion.  At \(\delta=0\), Definition \(\texttt{def:clamp-policy}\) is the identity and Proposition \(\texttt{prop:pushforward-setup}\) gives \(q_x(0)=0\); at \(\delta=\bar\delta<1\), every evaluation remains in \([0,1]\).  This proves every conclusion over the full closed threshold range for the law's own packaged carrier.